\chapter{부품 자루}
% full version: chapters/01b_neurontypes.tex

\pencilfig{1}{\pencilart{1}}{부품 자루: 거의 모두 파란색이고, 빨간색이 조금, 나머지는 한 줌이다.}

겉보기에 똑같은 부품 860억 개가 든 자루를 받았다고 하자. 그리고 이 부품으로 컴퓨터가 어떻게 조립되어 있는지 알아내라는 과제를 받았다. 어디서부터 시작하겠는가? 엔지니어라면 부품의 모양을 뜯어보지 않을 것이다. 그는 이렇게 물을 것이다. 입력은 몇 개이고 출력은 몇 개이며, 출력으로 무엇을 내보내는가?

이런 부품을 뉴런이라고 부르는데, 구조는 놀라울 만큼 단순하다. 입력은 수천 개다. 다른 뉴런에서 온 전선이 여기에 닿는다. 출력은 하나다. 그러나 이 출력은 나무처럼 가지를 치고, 같은 신호의 복사본이 수천 명의 수신자에게 간다. 천 명의 말을 한꺼번에 듣고 한 문장으로 답하는데, 그 답을 또 다른 천 명이 듣는 사람을 떠올려 보라.

그 신호는 무엇일까? 짧은 딸깍이다. 뉴런은 이따금 스파이크를 낸다. 스파이크는 길이가 1000분의 1초 정도인 짧은 전압 펄스다. 한 뉴런의 딸깍은 한 북을 두드리는 소리처럼 모두 똑같다. 그래서 뉴런이 전할 수 있는 모든 것은 딸깍의 크기가 아니라 \emph{언제} 딸깍하는지에 담긴다. 더 자주인지 더 드물게인지, 곧바로인지 늦게인지.

뉴런 안에서 일어나는 계산은 상상할 수 있는 가장 단순한 것이다. 들어오는 신호를 더하고, 합이 문턱값을 넘으면 스스로 딸깍한다. 어떤 입력은 합을 위로 밀어 올리고, 어떤 입력은 아래로 끌어내린다.

\section*{자루를 분류하는 법}

출력 전선의 끝은 이웃 뉴런에 닿을 만큼 뻗어 가지만 실제로 맞닿지는 않는다. 둘 사이에는 좁은 틈이 남고, 신호는 화학 물질이 이 틈을 건너 전한다. 전선 끝이 특별한 물질, 곧 신경전달물질을 조금 뿜어내면 이웃이 그것을 받아낸다. 그런데 운 좋게도 거의 모든 뉴런은 전선의 모든 끝에서 같은 물질을 내보낸다. 그러니 뉴런을 \emph{무엇을} 내보내는지에 따라 분류할 수 있다.

이 책에서는 뉴런의 유형을 그 물질의 영어 이름 첫 네 글자로 부른다. 글루탐산(glutamate)을 내보내는 뉴런은 \nt{GLUT}이다. 감마아미노뷰티르산(GABA)을 내보내는 뉴런은 \nt{GABA}이다. 그다음으로 \nt{DOPA}(도파민), \nt{SERO}(세로토닌), \nt{NORA}(노르아드레날린), \nt{ACET}(아세틸콜린) 등 몇 가지가 더 있다.

자루를 다 분류하고 나면 분포가 몹시 고르지 않다. 뉴런 천 개 가운데 960개 정도가 \nt{GLUT}이다. 이들의 신호는 수신자를 딸깍 쪽으로 밀어 준다. 이것이 “플러스”다. 36개 정도는 \nt{GABA}이고, 이들의 신호는 수신자를 문턱값에서 멀어지게 한다. 이것이 “마이너스”다. 대뇌 피질에서는 GABA 뉴런이 더 많아서 5분의 1에서 4분의 1에 이른다. 나머지 유형은 모두 합쳐도 10만 개 가운데 몇 개에 불과하다. 바다의 물 한 방울이다.

\section*{전화와 라디오}

그러나 이 한 방울은 전혀 다르게 작동한다. \nt{GLUT} 뉴런이나 \nt{GABA} 뉴런의 전선은 전화선처럼 정해진 이웃으로 이어진다. 전화를 건 상대만 신호를 받는다. 반면 \nt{DOPA}, \nt{SERO}, \nt{NORA}, \nt{ACET}의 작은 무리는 방송국처럼 일한다. 이들의 전선은 뇌의 넓은 영역으로 퍼져 나가고, 물질은 틈이 아니라 주변 공간으로 흘러나오는 일이 많다. 그래서 같은 메시지를 수백만 개의 세포가 한꺼번에 듣는다.

전화선으로는 데이터가 흐른다. 그림의 이 선, 이 높이의 소리 같은 것이다. 라디오로는 전체 설정이 흐른다. 모든 것을 얼마나 크게 할지, 지금 학습할지, 잘 시간인지. 어느 컴퓨터나 마찬가지다. 메모리 셀은 수십억 개이지만 제어 레지스터는 한 줌이다. 그래도 제어 레지스터가 없으면 기계는 돌아가지 않는다.

\section*{의미는 수신자가 정한다}

한 가지 미묘한 점이 있다. 물질 자체는 자신이 “플러스”인지 “마이너스”인지 모른다. 분자는 열쇠다. 열쇠를 꽂았을 때 무슨 일이 일어날지는 자물쇠가 정한다. 자물쇠는 수신자 쪽에 있는 특별한 단백질, 곧 수용체다. 같은 도파민이 어떤 수신자는 북돋우고 어떤 수신자는 누그러뜨린다. 어떤 자물쇠가 달려 있느냐에 따라 다르다. 같은 라디오 방송을 청취자마다 다르게 알아듣는 것과 같다.

\section*{“예상보다 좋다”는 신호}

가장 유명한 라디오 채널은 도파민이다. 도파민은 무엇을 전할까? 실험을 하나 보자. 동물에게 예고 없이 주스 한 방울을 주면 도파민 뉴런이 딸깍을 한바탕 쏟아낸다. 다음에는 주스를 주기 전에 매번 전구를 켠다. 동물은 전구가 주스를 약속한다는 것을 배우고, 이제는 전구에 딸깍이 쏟아지고 주스 자체에는 쏟아지지 않는다. 전구를 켜고 주스를 주지 않으면, 도파민 뉴런은 주스를 기대하던 순간에 \emph{조용해진다}.

\pencilfig{101}{%
% three rows: a spike raster of a DOPA neuron; lamp (amber) and juice (blue drop) above the axis
\foreach \row/\y in {1/0,2/-1.05,3/-2.1}{
  \draw[pen] (0,\y) -- (6.2,\y);
  \foreach \s in {0.3,0.75,1.1,2.15,2.6,3.05,3.45,3.9,4.45,4.95,5.4,5.85}{
    \pgfmathtruncatemacro{\gap}{(\row==3)&&(\s>3.6)&&(\s<4.7)}
    \ifnum\gap=0 \draw[pcGray, line width=0.7pt] (\s,\y) -- (\s,\y+0.3);\fi}}
% row 1: juice without a lamp -> burst at the juice
\foreach \s in {4.0,4.08,4.16,4.24,4.33}{\draw[pcAmber, line width=0.8pt] (\s,0) -- (\s,0.45);}
\draw[dotpen=pcBlue, wash=pcBlue] (4.15,0.72) circle (0.09); \draw[pcBlue, line width=0.6pt] (4.07,0.76) -- (4.15,0.92) -- (4.23,0.76);
% row 2: lamp -> burst at the lamp, nothing at the promised juice
\foreach \y in {-1.05,-2.1}{
  \foreach \s in {1.5,1.58,1.66,1.74,1.83}{\draw[pcAmber, line width=0.8pt] (\s,\y) -- (\s,\y+0.45);}
  \draw[dotpen=pcAmber, wash=pcAmber] (1.65,\y+0.72) circle (0.1);
  \foreach \a in {30,90,150}{\draw[pcAmber, line width=0.5pt] ($(1.65,\y+0.72)+(\a:0.14)$) -- ($(1.65,\y+0.72)+(\a:0.24)$);}}
\draw[dotpen=pcBlue, wash=pcBlue] (4.15,-0.33) circle (0.09); \draw[pcBlue, line width=0.6pt] (4.07,-0.29) -- (4.15,-0.13) -- (4.23,-0.29);
% row 3: lamp, no juice -> a gap where the juice was expected
\draw[pcBlue, line width=0.6pt, densely dotted] (4.15,-1.38) circle (0.09); \draw[pcBlue, line width=0.6pt, densely dotted] (4.07,-1.34) -- (4.15,-1.18) -- (4.23,-1.34);
\draw[penlite=pcRed] (3.65,-2.22) -- (3.65,-2.28) -- (4.65,-2.28) -- (4.65,-2.22);
\node[lbl, text=pcRed, anchor=north] at (4.15,-2.3) {침묵};
\node[lbl, anchor=east, align=right] at (-0.1,0.15) {전구 없이\\주스};
\node[lbl, anchor=east, align=right] at (-0.1,-0.9) {전구,\\다음에 주스};
\node[lbl, anchor=east, align=right] at (-0.1,-1.95) {전구,\\주스 없음};
}{예고 없는 주스에 쏟아지는 딸깍, 주스를 약속하는 전구에 쏟아지는 딸깍, 그리고 약속한 주스가 오지 않은 순간의 침묵.}

세 경우를 모두 설명하는 규칙은 이것이다. 도파민이 알리는 것은 “좋다”가 아니라 “예상보다 좋다”이다. 예고 없는 주스는 반가운 깜짝 선물이다. 약속된 주스는 아무런 놀라움도 아니다. 약속했는데 주지 않은 주스는 불쾌한 놀라움이다. 자신의 실수에서 배우는 프로그램이 쓰는 값이 바로 이것이다. 이 값은 나중에 다시 다룬다.

\section*{기억해 둘 것}

뇌는 겉보기에 똑같은 부품으로 이루어진 거대한 네트워크다. 그 부품의 거의 전부는 데이터를 나르는 전화망의 “플러스”와 “마이너스”다. 그리고 아주 적은 수가 네트워크 전체를 한꺼번에 조율하는 방송국이다. 다음 장에서는 “플러스”만으로 무엇을 만들 수 있는지 알아본다.

\fullversion{1}
